\chapter{ПЕРЕЛІК УМОВНИХ ПОЗНАЧЕНЬ, СИМВОЛІВ, ОДИНИЦЬ, СКОРОЧЕНЬ І ТЕРМІНІВ}
% За алфавітом


\begin{enumerate}
  \item Вступ.
  \item Огляд літератури
  \item Формулювання задач і методів.
    \begin{enumerate}
      \item Плоска задача теплопровідності у тілі
      \item Плоска задача пружності
      \item Плоска задача термопружності
      \item Метод граничних елемен­тів
      \item Метод пригранич­них елемен­тів
    \end{enumerate}
  \item Плоска задача теплопровідності у тілі
    \begin{enumerate}
      \item Формулювання задачі із круглим включенням
      \item Дискретизація рівнянь 
      \item Розв'язок
      \item Формулювання задачі із багатьма включеннями
      \item Дискретизація рівнянь 
      \item Розв'язок
      \item Висновки
    \end{enumerate}
  \item Великий шрам у місці розтину.
\end{enumerate}


\begin{description}
  \item[КМГЕ] --- комплексний метод граничних елементів 
  \item[МГЕ] --- метод граничних елементів
  \item[МПГЕ] --- метод приграничних елементів
  \item[МСЕ] --- метод скінчених елементів
  \item[МФО] --- метод фіктивних областей
  \item[РМСЕ] --- розширений метод скінчених елементів
  \item[МСС] --- метод скінчених секцій 
  \item[МСЕФО] --- метод скінчених елементів фантомних областей 
  \item[ГМТСЕ] --- гібридний метод Треффца і скінчених елементів 
\end{description}